\lecture{25}
\title{Logistic Regression}

\begin{document}

We now switch from linear regression to logistic regression.
The only difference is that the response variable $Y$
now takes on a binary value: $Y \in \{0, 1\}$.

\textbf{Definition.}
Given data $\{(X_i, Y_i)\}_{i = 1}^n$,
with $X_i \in \mathbb{R}^k$ and $Y_i \in \{0, 1\}$,
our \textbf{logistic regression} model will be:
\[ [Y \mid X = x] \sim \mathrm{Ber}\left(
    \sigma(x^{\top}\beta^*)
\right) ~~~ \text{ where } ~~~
\sigma(t) := \frac{e^t}{1 + e^t}. \]
Note that $\sigma(t)$ is the \textbf{sigmoid function},
pictured below.
\begin{center}
    \frame{\includegraphics[width=0.4\textwidth]{lecture-25/sigmoid}}
\end{center}
Note that the output $\sigma(x^{\top}\beta^*)$
of our regression model
is not an estimate of the response,
but rather an estimate of the probability
that the response is equal to $1$.

\begin{remark}
    Another commonly-used model is the
    \textbf{probit regression model},
    described below:
    \[ [Y \mid X = x] \sim \mathrm{Ber}\left(
        \Phi\left(x^{\top}\beta^*\right)
    \right) ~~~ \text{ where } ~~~
    \Phi(t) \text{ is the CDF of } \mathcal{N}(0, 1). \]
    An equivalent representation of the above model is:
    \[ Y = \mathds{1}(X^{\top}\beta^* + Z > 0)
    ~~~ \text{ where } ~~~
    Z \sim \mathcal{N}(0, 1). \]
    Not much is special about $\mathcal{N}(0, 1)$, either.
    Any function $F(t)$ can replace the sigmoid function $\sigma(t)$,
    so long as:
    \[ F(t) \text{ is increasing.}
    ~~~ \parallel ~~~
    \lim_{t \to -\infty} F(t) = 0 \text{ and }
    \lim_{t \to \infty} F(t) = 1.
    ~~~ \parallel ~~~
    F(0) = \frac{1}{2}. \]
\end{remark}

\subsection{MLE for Logistic Regression}

Unfortunately, there is no closed-form expression
for the MLE $\hat{\beta}^{\mathrm{MLE}}$ of logistic regression.
But we can still do our best
to try to approximate the log-likelihood-maximizer
using gradient ascent.

\begin{theorem}{Logistic Likelihood}
    The log-likelihood function for logistic regression is:
    \[ \ell_n(\beta)
    = \sum_{i = 1}^n \left[
        Y_iX_i^{\top}\beta - \log\left(
            1 + e^{X_i^{\top}\beta}
        \right)
    \right]. \]
\end{theorem}

\begin{proof}
    Recalling that $Y_i \mid X_i \sim \mathrm{Ber}\left(
        \sigma\left(X_i^{\top}\beta\right)
    \right)$, the log likelihood is just
    \begin{align*}
        \ell_n(\beta)
        & = \sum_{i = 1}^n \log\left(
            \sigma(X_i^{\top}\beta)^{Y_i} \cdot
            \left(1 - \sigma(X_i^{\top}\beta)\right)^{1 - Y_i}
        \right) \\
        & = \sum_{i = 1}^n \left[
            Y_i \log\left(
                \frac{\sigma(X_i^{\top}\beta)}{1 - \sigma(X_i^{\top}\beta)}
            \right)
            + \log(1 - \sigma(X_i^{\top}\beta))
        \right] \\
        & = \sum_{i = 1}^n \left[
        Y_iX_i^{\top}\beta - \log\left(
            1 + e^{X_i^{\top}\beta}
        \right)
        \right] ~~~ \text{ using the identities }
        \frac{\sigma(t)}{1 - \sigma(t)} = e^t
        \text{ and }
        1 - \sigma(t) = \frac{1}{1 + e^t}.
    \end{align*}
    And that's the furthest simplification can take us. ~ $\blacksquare$
\end{proof}

The upside is that
the log likelihood is concave (proof omitted),
so plain gradient ascent should work.

\subsection{Multiclass Classification}

A natural generalization of a binary response $Y \in \{0, 1\}$
is a multiclass response $Y \in \{0, 1, \dots, M\}$.

The result of a multiclass logistic regression
is a collection of $M$ parameters
$(\beta_1^*, \beta_2^*, \dots, \beta_M^*)$.
These $M$ parameters collectively decide the behavior
of $M + 1$ probability functions $(p_0(x), p_1(x), \dots, p_M(x))$,
where $\mathbb{P}[Y = m \mid X = x] = p_m(x)$.

\begin{remark}
    Check that this description is consistent with
    binary-classification logistic regression when $M = 1$.

    Also note that there are $M$ parameters,
    yet $M + 1$ probability functions; this is because
    the probability functions are constrained by the requirement
    that $\sum_{m = 0}^M p_m(x) = 1$ for all $x$.
\end{remark}

So how exactly do the parameters $\{\beta_m^*\}_{m = 1}^M$
relate to the probabilities $\{p_m(x)\}_{m = 0}^M$
in multiclass logistic regression?

\textbf{Definition.}
Given data $\{(X_i, Y_i)\}_{i = 1}^n$,
with $X_i \in \mathbb{R}^k$ and $Y_i \in \{0, 1, \dots, M\}$,
our \textbf{multiclass logistic regression} model will be:
\[ \mathbb{P}[Y = m \mid X = x] = p_m(x)
~~~ \parallel ~~~
p_0(x) = \frac{1}{1 + \sum_{j = 1}^M e^{x^{\top}\beta_j^*}}
~~~ \text{ and } ~~~
p_m(x) = \frac{e^{x^{\top}\beta_m^*}}{
    1 + \sum_{j = 1}^M e^{x^{\top}\beta_j^*}
} ~ \text{ for all } m \neq 0. \]

\begin{remark}
    It turns out that the log likelihood in multiclass logistic regression
    is the negative of the cross entropy loss from machine learning.
    (18.650 does not elaborate on this any further.)
\end{remark}

\end{document}